\chapter{Funzione Hamiltoniana}
\section{La trasformata di Legendre}

Sia \(f:\mathbb{R}\longrightarrow\mathbb{R}\) almeno derivabile con continuità due volte. Assumiamo inoltre che
\[
f''(x)>0 \quad \forall x\in\mathbb{R}.
\]
Questa condizione garantisce che la derivata prima
\[
f'(x) \text{ sia strettamente monotona crescente}.
\]

Possiamo allora invertire \(f'(x)\), ovvero definire
\[
x=(f')^{-1}(p) \quad \forall p\in \operatorname{Im}(f').
\]

Una possibile costruzione per la trasformata di Legendre è quella geometrica.

Fisso una pendenza \(p\) e disegno la retta \(y=px\).

\begin{center}
\begin{tikzpicture}[scale=1.1, >=Stealth]
  % assi
  \draw[->, thick] (-2.2,0) -- (3.2,0);
  \draw[->, thick] (0,-1.5) -- (0,3.2);

  % parametro della retta
  \def\p{1.3}

  % funzione convessa di esempio
  \draw[thick, domain=-0.8:2.7, samples=200]
    plot (\x,{0.6*(\x-1)^2 + 0.4});
  \node at (-0.95,2.35) {\(f\)};

  % retta y = px
  \draw[thick] (-1.1,{-1.1*\p}) -- (2.35,{2.35*\p});
  \node at (1.3,2.75) {\(y=px\)};

  % punto della funzione in cui la tangente è parallela alla retta
  % f'(x)=1.2(x-1), quindi f'(x)=p per x=1+p/1.2
  \pgfmathsetmacro{\xt}{1+\p/1.2}
  \pgfmathsetmacro{\yf}{0.6*(\xt-1)^2 + 0.4}
  \pgfmathsetmacro{\yl}{\p*\xt}

  % segmento verticale rosso tra funzione e retta
  \draw[red, thick] (\xt,\yf) -- (\xt,\yl);

  % freccia e scritta rossa
  \draw[red, thick, ->] (\xt+0.08,{(\yf+\yl)/2})
    .. controls (\xt+0.55,{(\yf+\yl)/2+0.2}) and (\xt+0.95,{(\yf+\yl)/2+0.1}) ..
    (\xt+1.15,{(\yf+\yl)/2+0.15});
  \node[red] at (\xt+1.55,{(\yf+\yl)/2+0.35}) {\(f^*(p)\)};
\end{tikzpicture}
\end{center}

Definisco
\[
f^*(p)=xp-f(x)
\]
con \(x\) tale che \(px-f(x)\) è massima ovvero
\[
f^*(p)=\sup_{x\in\mathbb{R}}\bigl(xp-f(x)\bigr).
\]

Tale valore è assunto per \(x\in\mathbb{R}\) tale che \(f'(x)=p\), ovvero
\[
x(p)=(f')^{-1}(p).
\]

Infatti dobbiamo massimizzare
\[
g(x)=xp-f(x)
\quad \Rightarrow \quad
g'(x)=p-f'(x)
\]
\[
\Rightarrow \quad
g''(x)=-f''(x).
\]
Allora
\[
g'(x)=0 \Longleftrightarrow f'(x)=p
\]
ed inoltre
\[
g''(x)<0 \quad \Rightarrow \quad x \text{ è un massimo}.
\]

Spesso la trasformata di Legendre viene quindi indicata come
\[
f^*(p)=x(p)p-f(x(p)),
\qquad
x(p)=(f')^{-1}(p).
\]

\begin{proposition}{}
Sia \(f(x)\) strettamente convessa su \(\mathbb{R}\) e almeno \(C^2(\mathbb{R})\) e sia \(f^*(p)\) la sua trasformata di Legendre, allora
\[
\frac{d f^{(*)}(p)}{dp}=x
\]
ovvero l'antitrasformata della trasformata di Legendre è la sua derivata.
\end{proposition}

\begin{proof}
\[
\frac{d f^{(*)}(p)}{dp}
=
\frac{d}{dp}\bigl(xp-f(x)\bigr)
=
\frac{dx}{dp}p+x-f'(x)\frac{dx}{dp}
=
\frac{dx}{dp}p+x-p\frac{dx}{dp}
=
x.
\]
\end{proof}

\section{Funzione di Hamilton}

Consideriamo \(\mathcal{L}(q_1,\ldots,q_m,\dot q_1,\ldots,\dot q_m)\) lagrangiana per un sistema a \(m\) gradi di libertà.
Sappiamo che possiamo scrivere le equazioni di Lagrange
\[
\frac{d}{dt}\frac{\partial \mathcal{L}}{\partial \dot q_j}
=
\frac{\partial \mathcal{L}}{\partial q_j}
\qquad
j=1,\ldots,m.
\]

Sappiamo anche che se
\[
\frac{\partial \mathcal{L}}{\partial q_j}=0
\quad \Rightarrow \quad
\frac{\partial \mathcal{L}}{\partial \dot q_j}
\]
è un integrale primo.

In generale un moto descritto nel piano delle fasi è una curva regolare

\begin{center}
\begin{tikzpicture}[scale=1.15, >=Stealth]
  \draw[->, thick] (0,0) -- (0,2.4);
  \draw[->, thick] (0,0) -- (2.4,0);
  \draw[thick] (-1.0,-1.2) -- (0,0);

  \draw[thick]
    plot[smooth, tension=0.8] coordinates {
      (-1.3,-0.8)
      (-0.9,0.0)
      (-0.25,0.55)
      (0.35,0.7)
      (0.95,0.65)
      (1.45,0.78)
      (1.9,1.05)
    };

  \node at (2.9,1.15) {\((q(t),\dot q(t))\)};
\end{tikzpicture}
\end{center}

Se però riusciamo ad usare
\[
\frac{\partial \mathcal{L}}{\partial \dot q_j}
\]
come nuova coordinata libera \(\tilde q\)
allora il moto sarebbe descritto da una curva piana

\begin{center}
\begin{tikzpicture}[scale=1.15, >=Stealth]
  \draw[->, thick] (0,0) -- (0,2.5);
  \draw[->, thick] (0,0) -- (2.5,0);
  \draw[->, thick] (0,0) -- (-1.2,-1.25);

  \node at (-0.25,2.65) {\(\tilde q\)};

  \draw[red, thick]
    plot[smooth, tension=0.9] coordinates {
      (-0.85,-1.0)
      (-0.45,-0.45)
      (0.25,-0.18)
      (0.9,0.08)
      (1.4,0.28)
      (1.65,0.55)
      (1.95,0.48)
    };

  \node at (-0.45,-1.35) {\((q(t),\dot q(t))\)};

  \draw[thick, ->]
    (1.65,0.25) .. controls (2.0,-0.05) and (2.4,0.0) .. (2.65,0.1);

  \node at (3.35,0.35) {sta nel piano};
  \node at (3.15,-0.15) {\(\dot{\tilde q}=\text{cost.}\)};
\end{tikzpicture}
\end{center}

ovvero lavoreremmo in uno spazio delle fasi di dimensione minore.

Questa è l'idea dietro la funzione di Hamilton. Come vedremo però la sua definizione non dipende da integrali primi del moto ed è quindi utilizzabile anche in casi più generali.

\begin{definition}{}
Si definisce funzione Hamiltoniana la trasformata di Legendre della lagrangiana.
\end{definition}

Se consideriamo un sistema ad \(1\) grado di libertà siamo già pronti a scrivere l'Hamiltoniana.
\[
\mathcal{L}=\mathcal{L}(q,\dot q)
\qquad
\text{definisco}
\qquad
p=\frac{\partial \mathcal{L}}{\partial \dot q}
\qquad
\text{allora}
\]
\[
H(p,q)=p\dot q(p,q)-\tilde{\mathcal{L}}(q,p)
\]
avendo definito
\[
\tilde{\mathcal{L}}(q,p)=\mathcal{L}(q,\dot q(q,p)).
\]

Possiamo generalizzare la definizione al caso con \(m\) gradi di libertà ma prima facciamo un'osservazione \((\text{valida anche per }m=1)\).
\begin{oss}{}
\(\mathcal{L}\) è una funzione convessa di \(\dot q\).
\[
\mathcal{L}
=
T(\underline q,\dot{\underline q})-V(\underline q)
=
\frac{1}{2}\sum_{i=1}^{n}m_i\,\underline v_i\cdot \underline v_i
-
V(\underline q)
=
\frac{1}{2}\sum_{i=1}^{n}m_i\,
\frac{d\underline x_i}{dt}\cdot \frac{d\underline x_i}{dt}
-
V(\underline q)
=
\]
\[
=
\frac{1}{2}\sum_{i=1}^{n}
\left[
m_i
\left(
\sum_{j=1}^{m}\frac{\partial \underline x_i}{\partial q_j}\dot q_j
\right)
\left(
\sum_{k=1}^{m}\frac{\partial \underline x_i}{\partial q_k}\dot q_k
\right)
\right]
-
V(\underline q)
=
\frac{1}{2}\sum_{j=1}^{m}\sum_{k=1}^{m}
\left(
\sum_{i=1}^{n}m_i
\frac{\partial \underline x_i}{\partial q_j}\cdot
\frac{\partial \underline x_i}{\partial q_k}
\right)
\dot q_j\dot q_k -V(\underline{q})
=
\]
\[
=
\frac{1}{2}\sum_{j=1}^{m}\sum_{k=1}^{m}
a_{jk}\dot q_j\dot q_k
-
V(\underline q)
=
\frac{1}{2}\dot{\underline q}^{T}\underline a\,\dot{\underline q}
-
V(\underline q).
\]

Dal momento che
\[
T=
\frac{1}{2}\dot{\underline q}^{T}\underline a\,\dot{\underline q}\ge 0
\qquad
\text{e}
\qquad
T=0 \Longleftrightarrow \dot{\underline q}^{T}=\underline 0
\]
allora \(T\) è una forma quadratica definita positiva e \(\mathcal{L}(\underline q,\dot{\underline q})\) è una funzione convessa.

Nel caso \(m\) dimensionale la buona definizione dell'Hamiltoniana non dipende tanto dalla convessità di \(\mathcal{L}\) ma più dal fatto che quest'ultima è una forma quadratica definita positiva.

Definiamo
\[
p_j=
\frac{\partial \mathcal{L}}{\partial \dot q_j}
=
\frac{\partial T}{\partial \dot q_j}
\]
perché \(V(q)\) non dipende da \(\dot q_j\). Allora
\[
\underline p
=
(p_j)_{j=1,\ldots,m}
=
\nabla_{\dot{\underline q}}
\left(
\frac{1}{2}\dot{\underline q}^{T}\underline a\,\dot{\underline q}
\right)
=
\underline a\,\dot{\underline q}.
\]

Ovvero la mappa
\[
\dot{\underline q}\longmapsto \underline p
\]
è un'applicazione lineare di matrice \(\underline a\).
A questo punto usiamo il fatto che \(T\) è definita positiva e quindi ha autovalori reali positivi che garantiscono l'invertibilità.
\end{oss}
Per un sistema a \(m\) coordinate libere possiamo scrivere
\[
\mathcal{L}
=
\mathcal{L}(q_1,\ldots,q_m,\dot q_1,\ldots,\dot q_m)
=
\mathcal{L}(\underline q,\dot{\underline q})
\]
e
\[
H(\underline q,\underline p)
=
\underline p\cdot \dot{\underline q}(\underline q,\underline p)
-
\tilde{\mathcal{L}}(\underline q,\underline p)
=
\sum_{j=1}^{m}p_j\dot q_j
-
\tilde{\mathcal{L}}(q_1,\ldots,q_m,p_1,\ldots,p_m)
\]
con \(\underline p\) definito sopra e
\[
\tilde{\mathcal{L}}(\underline q,\underline p)
=
\mathcal{L}(\underline q,\dot{\underline q}(\underline q,\underline p)).
\]

Calcoliamo ora il differenziale di \(H(q,p)\) che ricordiamo agisce sullo spazio tangente.
Partiamo dal caso \(m=1\)
\[
\nabla H(q,p)
\binom{\delta q}{\delta p}
=
\begin{pmatrix}
\dfrac{\partial H}{\partial q} & \dfrac{\partial H}{\partial p}
\end{pmatrix}^{T}
\cdot
\binom{\delta q}{\delta p}
=
-\frac{\partial}{\partial q}\bigl(\tilde{\mathcal{L}}(q,p)\bigr)\cdot\delta q
+
\frac{\partial}{\partial q}\bigl(p\dot q(p,q)\bigr)\cdot\delta q
+
\]
\[
+
\frac{\partial}{\partial p}\bigl(p\dot q(p,q)\bigr)\delta p
-
\frac{\partial}{\partial p}\bigl(\tilde{\mathcal{L}}(q,p)\bigr)\delta p
=
\]
\[
=
-\frac{\partial}{\partial q}\bigl(\mathcal{L}(q,\dot q(q,p))\bigr)\cdot\delta q
+
p\frac{\partial \dot q}{\partial q}\cdot\delta q
+
\dot q(p,q)\delta p
+
p\frac{\partial \dot q}{\partial p}\delta p
-
\frac{\partial}{\partial p}\bigl(\mathcal{L}(q,\dot q(q,p))\bigr)\cdot\delta p
=
\]
\[
=
-\frac{\partial\mathcal{L}}{\partial q}\delta q
-
\frac{\partial\mathcal{L}}{\partial \dot q}\cdot
\frac{\partial\dot q}{\partial q}\cdot\delta q
+
p\frac{\partial\dot q}{\partial q}\cdot\delta q
+
\dot q\cdot\delta p
+
p\frac{\partial\dot q}{\partial p}\delta p
-
\frac{\partial\mathcal{L}}{\partial \dot q}\cdot
\frac{\partial\dot q}{\partial p}\cdot\delta p
\]

Ricordando che
\[
\frac{\partial\mathcal{L}}{\partial \dot q}=p
\]
\[
\nabla H(q,p)
\binom{\delta q}{\delta p}
=
-\frac{\partial\mathcal{L}}{\partial q}\cdot\delta q
-
p\cdot\frac{\partial\dot q}{\partial q}\cdot\delta q
+
p\frac{\partial\dot q}{\partial q}\cdot\delta q
+
\dot q\delta p
+
p\frac{\partial\dot q}{\partial p}\delta p
-
p\cdot\frac{\partial\dot q}{\partial p}\delta p
=
\]
\[
=
-\frac{\partial\mathcal{L}}{\partial q}\delta q
+
\dot q\delta p
\]

Uso l'equazione di Lagrange
\[
\frac{d}{dt}\frac{\partial\mathcal{L}}{\partial \dot q}
=
\dot p
=
\frac{\partial\mathcal{L}}{\partial q}
\]
\[
\nabla H(q,p)
=
-\dot p\delta q+\dot q\delta p
\]

Dal momento che è valido \(\forall \delta q,\delta p\) allora ottengo le equazioni di Hamilton
\[
\begin{cases}
\dot q=\dfrac{\partial H}{\partial p}\\[6pt]
\dot p=-\dfrac{\partial H}{\partial q}
\end{cases}
\]


Nel caso a \(m\) gradi di libertà con passaggi analoghi dall'Hamiltoniana
\[
H(\underline q,\underline p)
=
\underline p\cdot \dot{\underline q}(\underline q,\underline p)
-
\tilde{\mathcal{L}}(\underline q,\underline p)
\]
si ottengono le equazioni di Hamilton
\[
\begin{cases}
\dot q_j=\dfrac{\partial H}{\partial p_j}\\[6pt]
\dot p_j=-\dfrac{\partial H}{\partial q_j}
\end{cases}
\qquad
\forall j=1,\ldots,m.
\]

Rispetto alle equazioni di Lagrange siamo passati da \(m\) equazioni del secondo ordine a \(2m\) equazioni del primo ordine.

\begin{oss}{}
\(m=1\)
\[
\begin{cases}
\dot q=\dfrac{\partial H}{\partial p}\\[6pt]
\dot p=-\dfrac{\partial H}{\partial q}
\end{cases}
\qquad
\binom{\dot q}{\dot p}
=
\begin{pmatrix}
0 & 1\\
-1 & 0
\end{pmatrix}
\binom{\dfrac{\partial H}{\partial q}}{\dfrac{\partial H}{\partial p}}
\]

La matrice
\[
\begin{pmatrix}
0 & 1\\
-1 & 0
\end{pmatrix}
\]
è detta matrice simplettica.
\end{oss}

\begin{oss}{}
Se
\[
\mathcal{L}=\mathcal{L}(\underline q,\dot{\underline q}_j,t)
\]
dipende anche dal tempo allora anche
\[
H=H(\underline q,\underline p,t)
\]
dipenderà dal tempo e allo stesso modo si possono ricavare equazioni di Hamilton modificate.
\end{oss}

\begin{oss}{}
Se un sistema è conservativo \(H\) è l'energia totale \(H(q,p)\).

Dimostriamo il caso \(m=1\). Sappiamo che
\[
T=\frac{1}{2}a(q)\dot q^2,
\qquad
a>0,
\qquad
\mathcal{L}=\frac{1}{2}a(q)\dot q^2-V(q),
\qquad
p=\frac{\partial\mathcal{L}}{\partial\dot q}
=
a(q)\dot q.
\]
\[
\Rightarrow
H
=
p\dot q-\mathcal{L}(q,\dot q)
=
p\cdot\frac{p}{a(q)}
-
\mathcal{L}\left(q,\frac{p}{a(q)}\right)
=
p\cdot\frac{p}{a(q)}
-
\frac{1}{2}a(q)\dot q^2
+
V(q)
=
\frac{p^2}{a(q)}
-
\frac{1}{2}a(q)\frac{p^2}{a^2(q)}
+
V(q)
=
\]
\[
=
\frac{1}{2}\frac{p^2}{a(q)}
+
V(q)
=
\frac{1}{2}a(q)\dot q^2+V(q)
=
T+V.
\]
\end{oss}


\begin{example}{}
Oscillatore armonico
\[
m\ddot{x}=-Kx
\]
\[
T=\frac{m\dot{x}^2}{2}
\qquad
V=\frac{Kx^2}{2}
\qquad
q=x
\quad \Rightarrow \quad
\mathcal{L}=\frac{m\dot{x}^2}{2}-\frac{Kx^2}{2}
\]
\[
p=\frac{\partial \mathcal{L}}{\partial \dot{x}}=m\dot{x}
\]
\[
H(p,q)
=
\frac{m\dot{x}^2}{2}
+
\frac{Kx^2}{2}
=
\frac{mp^2}{2m^2}
+
\frac{Kx^2}{2}
=
\frac{p^2}{2m}
+
\frac{Kx^2}{2}
\]
\end{example}

\begin{example}{}
Consideriamo un punto materiale di massa \(m\) soggetto al vincolo
\[
y=\frac{x^2}{2}.
\]

\begin{center}
\begin{tikzpicture}[scale=1.1, >=Stealth]
  % assi
  \draw[thick] (-2.2,0) -- (2.8,0);
  \draw[thick] (0,-0.8) -- (0,2.5);

  % parabola
  \draw[thick, domain=-1.75:1.9, samples=200]
    plot (\x,{0.5*\x*\x});

  \node at (1.8,1.95) {\(y=\dfrac{x^2}{2}\)};

  % punto materiale
  \fill[black] (-1.35,{0.5*1.35*1.35}) circle (1.4pt);

  % forza peso
  \draw[red, thick, ->] (-1.35,{0.5*1.35*1.35}) -- (-1.35,0.15);
  \node[red] at (-1.65,0.65) {\(mg\)};
  \draw[red, thick] (-1.75,0.15) -- (-1.25,0.15);
\end{tikzpicture}
\end{center}

Scelgo come coordinata libera \(s(t)\)
\[
\begin{cases}
x=s(t)\\[4pt]
y=\dfrac{s^2(t)}{2}
\end{cases}
\Rightarrow
\begin{cases}
\dot{x}=\dot{s}\\[4pt]
\dot{y}=\dfrac{2s\dot{s}}{2}=s\dot{s}
\end{cases}
\]
\[
\Rightarrow
T
=
\frac{1}{2}m\dot{\underline{x}}\cdot\dot{\underline{x}}
=
\frac{1}{2}m(\dot{x}^2+\dot{y}^2)
=
\frac{1}{2}m(\dot{s}^2+s^2\dot{s}^2)
=
\frac{1}{2}m\dot{s}^2(1+s^2)
\]
\[
V=mgy=mg\frac{s^2}{2}
\Rightarrow
\mathcal{L}
=
m(1+s^2)\frac{\dot{s}^2}{2}
-
mg\frac{s^2}{2}.
\]

Costruiamo l'equazione di Lagrange.
\[
\frac{\partial\mathcal{L}}{\partial\dot{s}}
=
m(1+s^2)\dot{s}
\Rightarrow
\frac{d}{dt}\frac{\partial\mathcal{L}}{\partial\dot{s}}
=
m\ddot{s}(1+s^2)+m\dot{s}\cdot s\cdot 2\dot{s}
=
m\ddot{s}(1+s^2)+2ms\dot{s}^2
\]
\[
\frac{\partial\mathcal{L}}{\partial s}
=
m\dot{s}^2s-mgs
\Rightarrow
m(1+s^2)\ddot{s}+s\dot{s}^2m=-mgs
\Rightarrow
(1+s^2)\ddot{s}+s\dot{s}^2=-gs
\]

Come sempre siamo arrivati ad un'equazione pura del moto del secondo ordine.
Proviamo l'approccio con le equazioni di Hamilton
\[
p=
\frac{\partial\mathcal{L}}{\partial\dot{s}}
=
m(1+s^2)\dot{s}
\Rightarrow
\dot{s}
=
\frac{p}{m(1+s^2)}
\]

In questo caso è stato possibile esplicitare un'equazione già da subito.
Questo non avviene in generale.
Calcoliamo l'hamiltoniana
\[
H
=
p\dot{s}-\mathcal{L}(s,\dot{s})
=
p\cdot\frac{p}{m(1+s^2)}
-
m\frac{(1+s^2)}{2}\cdot
\frac{p^2}{m^2(1+s^2)^2}
+
mg\frac{s^2}{2}
=
\frac{p^2}{2m(1+s^2)}
+
mg\frac{s^2}{2}
\]
\[
\Rightarrow
\begin{cases}
\dot{s}=\dfrac{\partial H}{\partial p}\\[6pt]
\dot{p}=-\dfrac{\partial H}{\partial s}
\end{cases}
\Rightarrow
\begin{cases}
\dot{s}=\dfrac{p}{(1+s^2)m}\\[8pt]
\dot{p}
=
-\left(
\dfrac{p^2(-2s)}{2m(1+s^2)^2}
+
mgs
\right)
=
\dfrac{p^2s}{m(1+s^2)^2}
-
mgs
\end{cases}
\]
Come previsto abbiamo ottenuto 2 equazioni del primo ordine accoppiate.
\end{example}
